% Claims worded to match docs/research/novelty-log.md and PRD §3.3 (revised 2026-10-14).
\section{Introduction}
\label{sec:intro}

A sensor node that cannot run the full model offloads its windows to a cloud model. The bytes it
sends are a choice: raw samples, a lossless or a bounded-lossy waveform, a few engineered
features, or a learned latent of 16 bytes. Smaller payloads meet deadlines on slow links and cost
less on metered ones. They can also change what the cloud decides.

Existing systems make this choice by heuristics or by learned policies, and none of them states
what the choice does to the downstream decision. AdaEdge~\cite{liu2024adaedge} already selects
among lossless and lossy codecs per segment by the agreement of the model's output on decoded data
with its output on raw data, but as a best-effort bandit with no bound. Conformal methods have
reached the edge~\cite{kalor2025blackbox,huang2025cascades,zhu2024cdcrc}, but they select models,
labels or resources, not the representation of the payload.

\paragraph{Certify, then optimise.} We separate the guarantee from the optimisation. A
sufficiency predictor scores, per window, how likely each representation is to change the cloud's
decision. Conformal risk control (CRC)~\cite{angelopoulos2024crc} calibrates one threshold on those
scores. The resulting safe set has the property that matters for a system: the bound on
certified disagreements holds for \emph{every} selector that picks inside the set
(Proposition~\ref{prop:anyselector}). The selector can use the link state, the queue or its own
randomness; it can be greedy, model-predictive or learned, and it can be replaced without
re-proving anything.

\paragraph{An honest evaluation.} The certificate is only as meaningful as the reference model it
agrees with, and the guarantee assumes exchangeable calibration and test data. We therefore
evaluate strictly on units never seen in training (bearings, people, speakers, sleepers), choose
every reference model by inner cross-validation on training units, and record every
candidate, losers included. On the bearing data, scores near 0.95 reported in the literature come
from splits that share bearings between training and test~\cite{realistic2026bearing}; under
leave-bearings-out splits our best reference model reaches a macro-F1 of 0.604. We pre-registered
seven hypotheses before the first evaluation and report every verdict as registered.

\paragraph{Findings.} The certificate is the result that holds. A cascade of a fixed set and a
per-window set, each at $\alpha/2$, keeps mean disagreement below $\alpha$ on every dataset we
reran with the final code; it was added after the first results and is exploratory. The
window-level certificate lands near $\alpha$, as a tight expected-risk bound predicts; it meets
its pre-registered threshold only on Sleep-EDF, not significantly, and the registered one-sided test
cannot reject for a bound that is tight at $\alpha$. Per-window adaptivity does not pay: it fails
its pre-registered value test on every dataset. On three of four an uncertified fixed choice is
cheaper, because when the certificate refuses the cheap latent the next certified payload is
kilobytes, which a slow link cannot deliver before the deadline; on HAR adaptive choice is cheaper
but misses the SLO-reduction target. A
mixed menu pays only where the menu holds a small representation close to the raw decision. These
are the regimes a practitioner needs to know.

\paragraph{Contributions.}
\begin{itemize}
  \item \textbf{C1 (narrowed).} A per-window sufficiency certificate over a heterogeneous
  representation menu, the any-selector statement (Proposition~\ref{prop:anyselector}), and a
  union-bound cascade that spends edge compute on per-window scoring only when the cost model says
  it could pay (Section~\ref{sec:method}).
  \item \textbf{C4.} A leakage-safe benchmark: leave-units-out splits, a recorded reference-model
  search, a discrete-event simulator validated against real sockets and a Docker testbed, and
  pre-registered hypotheses with their verdicts (Sections~\ref{sec:system}--\ref{sec:protocol}).
  \item \textbf{C5.} An empirical map of when adaptive representation choice pays, on four
  datasets rerun with the final code and three from earlier code, with negative results reported
  as such (Section~\ref{sec:results}).
\end{itemize}

We do not claim that task-aware compression, adaptive offloading, conformal prediction at the
edge or a heterogeneous codec menu is new; each has prior work (Section~\ref{sec:related}).
